\section{Glossary and Explainer Checklist}
\label{sec:glossary}

\subsection{Glossary}
\begin{table}[H]
\centering
\begin{tabular}{@{}p{0.24\linewidth}p{0.68\linewidth}@{}}
\toprule
Term & Definition \\
\midrule
PV & Premise Validation probability distribution over entail/contradict/neutral. \\
Active ($A$) & Evidence currently used by the reasoner. \\
Suspended ($S$) & Evidence held out unless recovery trigger opens it. \\
Counter ($C$) & Contradiction-oriented evidence retained for refutation signal. \\
Dual-stream & Two separate message-passing streams: support and refute. \\
Bridge rescue & Diffusion-based scoring to detect structurally useful edges. \\
Recovery & One-step move of selected edges from $S$ to $A$ under trigger gate. \\
Selective prediction & Ability to abstain based on uncertainty + insufficiency. \\
\AURC{} & Area under risk-coverage curve; lower is better. \\
Rationale faithfulness & Whether removing important edges reduces predicted-class confidence. \\
\bottomrule
\end{tabular}
\caption{Quick definitions for presentation and review settings.}
\end{table}

\subsection{Checklist for explaining Component 2 to others}
Use this sequence for clean technical explanation:
\begin{enumerate}
\item Start with the interface contract: claim row, triple row, and required probability constraints.
\item Explain why $A/S/C$ pools are separated and why Component 2 reasons on $A\cup C$ but scores bridges on $A\cup S\cup C$.
\item Show PV masks and dual-stream attention to clarify how support/refute signals are preserved.
\item Explain recovery gate: low \texttt{esi\_geom} is necessary but not sufficient; predicted bridge connectivity gain must be positive.
\item Present selective prediction and \AURC{} to show reliability beyond top-1 accuracy.
\item End with rationale faithfulness to show interpretability is tested, not only visualized.
\end{enumerate}

\subsection{Current limitations (for transparent academic reporting)}
\begin{itemize}
\item The reasoner is currently a deterministic NumPy implementation, not a trained deep model checkpoint.
\item Recovery is one-step and greedy top-$k$; no multi-round optimization is implemented.
\item Faithfulness criterion uses confidence-drop positivity; stronger causal tests can be added later.
\item M3 behavior depends on quality/completeness of Component 1 pair logs in compatibility mode.
\end{itemize}

